% q0063.tex — In Search of the Automatic Economy
\question{
  id        = {0063},
  title     = {In Search of the Automatic Economy},
  source    = {OBECON},
  year      = {2026},
  round     = {Seletiva IEO -- Phase C},
  language  = {English},
  topic     = {Micro},
  subtopic  = {Competitive Equilibrium / Time Scarcity},
  type      = {objective},
  statement = {%
    Consider a world where agents can build mechanisms that produce goods
    indefinitely without further material inputs, goods can be freely stored
    and exchanged, and there exists a medium of exchange with negligible
    intrinsic value. Which statement correctly describes competitive
    equilibrium in this setting?
  },
  choices   = {%
    \item Exchange is unnecessary; trade ceases in equilibrium.
    \item All relative prices equal one because marginal costs are zero.
    \item Relative prices reflect differences in the time required to obtain
          additional units.
    \item The price of a good is determined solely by the construction cost of
          the mechanism that produces it.
  },
  answer    = {C},
  solution  = {%
    In this economy, material resources are effectively unlimited, but time
    remains the binding scarce input. Goods that require more time per
    additional unit to obtain are relatively more expensive. Competitive
    equilibrium prices are therefore determined by marginal time requirements
    --- not by sunk mechanism-construction costs (which are irrelevant at the
    margin) and not by material costs (which are zero). Exchange remains
    valuable because agents have heterogeneous preferences and time endowments.
    (C) is correct.
  },
}
